\newpage
\section{Manual de Código y Despliegue}
\label{sec:manual_codigo}

El código fuente íntegro de la plataforma MoodNest, abarcando tanto la capa de presentación (\textit{frontend}) como la lógica de negocio (\textit{backend}) y los archivos de configuración de infraestructura, ha sido versionado, documentado y publicado de forma abierta para su libre consulta.

\subsection{Estructura del Repositorio}

El proyecto ha sido organizado siguiendo un enfoque de monorepositorio (\textit{Monorepo}), separando claramente los dominios tecnológicos en directorios independientes para facilitar su inspección, mantenimiento y escalabilidad:

\begin{itemize}
    \item \texttt{/backend}: Contiene el código fuente de la API REST desarrollada en Spring Boot (Java 17). Incluye la estructura clásica de \textit{Controllers}, \textit{Services}, \textit{Repositories} y \textit{Models}, así como la configuración de seguridad perimetral mediante Spring Security y JWT.
    \item \texttt{/frontend}: Alberga la Aplicación de Página Única (SPA) construida con React 18 y Vite. Se estructura internamente mediante la Arquitectura Basada en Componentes (\texttt{/components}, \texttt{/pages}, \texttt{/context}, \texttt{/services}).
    \item \texttt{docker-compose.yml}: Archivo de orquestación en la raíz del proyecto, encargado de definir y levantar los contenedores de los servicios (Backend, Frontend y base de datos MongoDB) de forma interconectada y automatizada.
\end{itemize}

\subsection{Acceso al Código Fuente}

Siguiendo los principios de transparencia, código abierto y buenas prácticas en la ingeniería de software, el código del proyecto se encuentra alojado de forma pública en la plataforma GitHub:

\begin{itemize}
    \item \textbf{Enlace al repositorio:} \url{https://github.com/guillermo-palacios/tfg-moodnest}
\end{itemize}

\subsection{Instrucciones de Compilación y Despliegue Local}

El proyecto ha sido contenerizado con Docker para garantizar un despliegue agnóstico e idéntico en cualquier sistema operativo, mitigando el problema clásico de la falta de dependencias en los entornos de desarrollo locales.

Para clonar la plataforma y levantarla por completo, el evaluador o desarrollador únicamente requiere tener instaladas las herramientas de \textbf{Git}, \textbf{Docker} y \textbf{Docker Compose} en su máquina anfitriona. Los pasos de ejecución son los siguientes:

\begin{enumerate}
    \item Clonar el repositorio público mediante la terminal de comandos:
    \begin{quote}
        \texttt{git clone https://github.com/guillermo-palacios/tfg-moodnest}
    \end{quote}
    \item Navegar hasta el directorio raíz del proyecto donde se ubica el archivo de orquestación:
    \begin{quote}
        \texttt{cd tfg-moodnest}
    \end{quote}
    \item Ejecutar el comando para construir las imágenes de Docker locales y levantar los servicios en segundo plano:
    \begin{quote}
        \texttt{docker compose up -d --build}
    \end{quote}
    \item Una vez finalizado el proceso de descarga y compilación, la plataforma estará completamente operativa y accesible a través de:
    \begin{itemize}
        \item \textbf{\url{http://localhost:5173}} 
    \end{itemize}
\end{enumerate}

Para detener de forma limpia el ecosistema de contenedores y liberar los puertos del sistema, se debe ejecutar el comando \texttt{docker compose down}.